% source: manti.01.ma, 1ère SM, devoir 1er semestre (Drive 0B6pfSJcEVng0c0o1ak4yX2p2TTg). Ex.3 logique / récurrence (hors thème)
\ex{Fonctions}{c0o1}{Ex. 1}{11}
On considère les fonctions $g$ et $h$ définies par $g(x)=\sqrt{x+2}$ et $h(x)=\dfrac{x}{x-1}$.
\begin{enumerate}
\item Dresser le tableau de variation de chacune des fonctions $g$ et $h$.
\item Vérifier que $g(2)=h(2)$ et tracer $C_g$ et $C_h$ dans un même repère.
\item Soit $f$ la fonction impaire définie sur $D=\mathbb{R}\setminus\{-1,1\}$ telle que :
$f(x)=h(x)$ si $x\in[0,2]\setminus\{1\}$ \ et \ $f(x)=g(x)$ si $x\geq2$.
\begin{enumerate}
\item Calculer $f(-3)$ et $f(-1)$ \ ([?] $-1\notin D$ dans l'original).
\item Calculer $f(x)$ pour $x$ appartenant à $]-\infty,-2]$.
\item Tracer la courbe de la fonction $f$ dans un autre repère.
\item Montrer que $(\forall x\in\mathbb{R})\ \big(x\in\,]1,2]\big)\Rightarrow\big(h(x)\geq2\big)$, puis étudier la monotonie de $f\circ f$ sur l'intervalle $]1,2]$.
\end{enumerate}
\end{enumerate}
\fin
\ex{Ensembles}{c0o1}{Ex. 2}{9}
\textbf{(I)} Soit $\alpha$ un nombre réel ; on considère l'ensemble $X_\alpha=\{x\in\mathbb{R}\ /\ |2x-\alpha|\leq4\}$.
\begin{enumerate}
\item Déterminer $\alpha$ pour que $]-1,1[\,\subset X_\alpha$.
\item Déterminer les valeurs de $\alpha$ pour lesquelles $X_\alpha\cap\,]-1,3[\,\neq\varnothing$.
\end{enumerate}
\textbf{(II)} Déterminer en extension les ensembles $A=\left\{x\in\mathbb{R}\ /\ \dfrac x2\geq\dfrac2x\right\}$ \ ([?] $A$ est un intervalle) et $B=\{x\in\mathbb{Z}\ /\ |x+2|+|x-2|=4\}$.

\textbf{(III)} On considère les ensembles $A=\left\{\dfrac{7\pi}{12}+\dfrac{k\pi}{4}\ /\ k\in\mathbb{Z}\right\}$ et $B=\left\{\dfrac\pi4+\dfrac{k'\pi}{6}\ /\ k'\in\mathbb{Z}\right\}$.
\begin{enumerate}
\item Vérifier que $\dfrac{5\pi}{6}\in A$. A-t-on $\dfrac{5\pi}{6}\in B$ ?
\item Déterminer l'intersection $A\cap B$.
\end{enumerate}
\fin
